%% P15 --- Funkcje wielu zmiennych i gradient
%%
%% Kazda liczba w tym programie, ktorej czytelnik nie policzy w pamieci,
%% pochodzi z code/p15_gradient.py i trafia na strone przez \val{}. Listing
%% konsolowy to zatwierdzony transkrypt zapisany przez ten skrypt.
%%
%% TEZA: KIERUNEK gradientu jest wyprowadzony, a nie wybrany. Pochodna
%% kierunkowa to iloczyn skalarny, wiec cosinus z Programu P05 rozstrzyga,
%% ktory kierunek jest najbardziej stromy, a ktore sa plaskie -- i "najbardziej
%% stromy" oraz "prostopadly do poziomicy" okazuja sie dwoma odczytami jednego
%% rownania.
%%
%% CO TEN PROGRAM DOSTAJE, odczytane z plikow, a nie zapamietane:
%%   F11  wyprowadza marsz w jednym wymiarze W CALOSCI -- w dol jest przeciwnie
%%        do znaku pochodnej, a x <- x - eta f'(x), przy czym minus nazwany
%%        jest tam "calym mechanizmem sterowania". Ten program NIE MOZE wiec
%%        wyprowadzac znaku ponownie. Oddaje za to te sama rekurencje na
%%        kierunek wlasny, co skrypt bramkuje.
%%   F12  mowi, ze d/dx nazywa, KTORA litera jest zmienna, i ze bedzie to
%%        "mialo ogromne znaczenie w Programie P15, gdzie jest ich wiecej".
%%   P05  daje a . b = |a| |b| cos(theta). Sekcja 3 to ta tozsamosc, w ktorej
%%        jednym z dwoch wektorow jest gradient. Nic wiecej nie jest potrzebne.
%%   P10  zbiera ksztalt niecki -- wartosci wlasne 20 i 1, elipsa 4,47 raza
%%        dluzsza niz szersza -- i mowi wprost, ze go nie wydaje. Sekcja 5
%%        wydaje go, na zatwierdzonych liczbach P10.
%%
%% CZEGO NIE RUSZA, sprawdzone wzgledem tools/programs.json:
%%   jakobiany i tryb odwrotny                                    -> P16
%%   hesjan, krzywizna i NAJWIEKSZY STABILNY KROK                  -> P17
%%   rachunek macierzowy, gradient softmax-cross-entropy           -> P18
%%   wypuklosc                                                     -> P19
%%   moment pedu, ktory NAPRAWIA zygzak produkowany tutaj          -> P20
%% Brief P20 mowi wprost, ze naprawia "zygzak z P15", wiec wyprodukowanie go i
%% zostawienie nienaprawionym jest kontraktem, a nie luka.
%%
%% SLOWO SIE ZDERZA i sekcja 2 to nazywa. F06 i F11 uzywaja slowa "gradient"
%% w znaczeniu WSPOLCZYNNIKA KIERUNKOWEGO prostej, co jest standardowe i bylo
%% tam potrzebne; tutaj jest to wektor. Takie samo potraktowanie P07 dal
%% slowu "wymiar" wykonujacemu trzy prace.
%%
%% Blizniak: programs/en/P15-gradient.tex. Ramka w ramke, makro w makro,
%% liczba w liczbe; tools/parity.py je porownuje.

\program{Funkcje wielu zmiennych i gradient}
\label{prog:P15}
\index{gradient}
\index{derivative!partial}

\begin{outcomes}
\outcome{1--7}{policzyć pochodną cząstkową, trzymając każde inne wejście w
  miejscu, i powiedzieć, na jakie pytanie ona odpowiada}
\outcome{8--13}{zebrać pochodne cząstkowe w gradient i odróżnić ten wektor od
  współczynnika kierunkowego prostej}
\outcome{14--19}{zapisać tempo zmiany wzdłuż dowolnego kierunku jako iloczyn
  skalarny z gradientem}
\outcome{20--28}{powiedzieć, dlaczego gradient jest kierunkiem najbardziej
  stromym i dlaczego spotyka poziomicę pod kątem prostym, z tego jednego
  iloczynu skalarnego}
\outcome{29--38}{powiedzieć, co marsz robi w długiej wąskiej dolinie i
  dlaczego}
\end{outcomes}

\begin{quiz}
\item $f(x, y) = x^{2}y + 3y$. Ile wynosi $\partial f/\partial x$?
  \teachesat{3--4}
  \answerto{$2xy$. Potraktuj $y$ jak stałą i różniczkuj tak, jak uczył
  Program~\ref{prog:F11} \dash{} składnik $3y$ nie ma w sobie $x$, więc nie
  wnosi nic.}
\item W którym kierunku, spośród wszystkich, $f$ rośnie najszybciej?
  \teachesat{20--21}
  \answerto{We własnym kierunku gradientu. Nie dlatego, że tak go
  zdefiniowano \dash{} dlatego, że tempo wzdłuż kierunku to gradient
  pomnożony skalarnie przez ten kierunek, a iloczyn skalarny jest największy,
  gdy oba wskazują tak samo.}
\item W którym kierunku porusza się spadek gradientu? \teachesat{29--30}
  \answerto{W kierunku gradientu \emph{ujemnego}. Gradient wskazuje drogę
  najszybszego wzrostu $f$, więc krok idzie przeciwnie do niego, a znak jest
  najczęstszym błędem w ręcznie pisanej pętli treningowej.}
\item Stoisz na poziomicy $f$ i idziesz wzdłuż niej. Jakie jest tempo zmiany
  $f$ wzdłuż twojej drogi? \teachesat{23--24}
  \answerto{Zero, z definicji poziomicy. A skoro to tempo jest gradientem
  pomnożonym skalarnie przez twój kierunek, gradient musi stać pod kątem
  prostym do poziomicy.}
\item Czy w długiej wąskiej dolinie ujemny gradient wskazuje najniższy punkt?
  \teachesat{31--32}
  \answerto{Nie, i o tym jest cała sekcja 5. Wskazuje w poprzek doliny, pod
  kątem prostym do poziomicy, na której stoisz, a to nie jest ten sam
  kierunek co droga do domu, chyba że niecka jest okrągła.}
\end{quiz}

% ============================================================
\section{Po jednym wejściu naraz}
\label{sec:P15-partials}
\index{derivative!partial}

\begin{fr}
Wszystko, co robiły Program~\ref{prog:F11} i Program~\ref{prog:F12}, dotyczyło
funkcji jednego wejścia. Strata taka nie jest. Jest funkcją każdego parametru
modelu naraz, a dla czegokolwiek wartego trenowania jest to liczba o dziewięciu
czy dziesięciu cyfrach.

Brzmi to jak inny przedmiot i nim nie jest. Program~\ref{prog:F12} zostawił
wskazówkę: powiedział, że $\frac{\mathrm{d}}{\mathrm{d}x}$ nazywa,
\emph{która} litera jest zmienną, i że nazwanie jej będzie miało znaczenie
tutaj.

Mając $f(x, y)$ i tylko maszynerię rachunku jednej zmiennej, jak w ogóle
wydobyłbyś z tego pochodną?
\dotline
\nextframe
\end{fr}

\begin{fr}
\begin{ansblock}
Trzymając jedno z nich w miejscu. Ustal $y$ na jakiejś wartości, a $f$ staje
się funkcją samego $x$, czyli czymś, co już umiesz różniczkować.
\end{ansblock}

To cała idea i nie ma w niej nic więcej. Zapisuje się to
$\frac{\partial f}{\partial x}$, z zaokrąglonym $\partial$ zamiast prostego
$\mathrm{d}$, i czyta \enquote{cząstkowa de $f$ po de $x$}. Zaokrąglona litera
przypomina, że coś innego jest trzymane w miejscu.

\begin{note}
Notacja nie niesie nowej matematyki. $\frac{\partial f}{\partial x}$ to zwykła
pochodna zwykłej funkcji jednej zmiennej \dash{} tej, którą dostajesz z $f$,
zamrażając każde wejście poza $x$. Każda reguła, którą dał
Program~\ref{prog:F12}, stosuje się bez zmian, bo nie ma tu nic, pod względem
czego miałyby być inne.
\end{note}

\mermaidfig{p15-hold-the-rest-still}{Zamrażanie pozostałych wejść i zbieranie
tego, co zostaje.}{trzymanie reszty}
\end{fr}

\begin{fr}
Spróbuj. $f(x, y) = x^{2}y + 3y$.

Zróżniczkuj względem $x$, traktując $y$ jak stałą.
\dotline
\nextframe
\end{fr}

\begin{fr}
\ans{$2xy$}
Pierwszy składnik to $y$ razy $x^{2}$, a $y$ jest tu stałą, więc wychodzi
przed: pochodna wynosi $y \cdot 2x$. Drugi składnik, $3y$, nie ma w sobie $x$
wcale, więc z punktu widzenia $x$ jest stałą i nie wnosi nic.

Teraz druga. $\frac{\partial f}{\partial y}$, przy $x$ trzymanym w miejscu.
\nextframe
\end{fr}

\begin{fr}
\ans{$x^{2} + 3$}
Tym razem to $x^{2}$ jest stałą mnożącą $y$, więc przeżywa jako ono samo, a $3y$
różniczkuje się do $3$.

\begin{trapbox}
Popatrz, co się właśnie stało z dwiema odpowiedziami.
$\frac{\partial f}{\partial x}$ nadal ma w sobie $y$, a
$\frac{\partial f}{\partial y}$ nadal ma $x$.

\textbf{Trzymanie wejścia w miejscu nie usuwa go z odpowiedzi.} Jest ono
trzymane w miejscu \emph{podczas różniczkowania}, a potem wraca do bycia
zmienną. Pochodna cząstkowa jest funkcją wszystkich tych samych wejść, co
oryginał \dash{} dlatego trzeba ją wyliczyć w punkcie, zanim będzie liczbą, i
dlatego \enquote{pochodna} sama w sobie nie jest jeszcze odpowiedzią na nic.
\end{trapbox}
\end{fr}

\begin{fr}
Wylicz je więc. W punkcie $x = \val{p15.pt.x}$, $y = \val{p15.pt.y}$ ile
wynoszą obie pochodne cząstkowe?
\dotline
\nextframe
\end{fr}

\begin{fr}
\ans{$\val{p15.fx.at}$ oraz $\val{p15.fy.at}$}
Z $2xy = 2 \times \val{p15.pt.x} \times \val{p15.pt.y}$ oraz
$x^{2} + 3 = \val{p15.pt.x}^{2} + 3$.

Obie zostały sprawdzone wobec różnicy centralnej na siatce punktów, tak jak
Program~\ref{prog:F12} sprawdzał swoje reguły, i zgadzają się wszędzie na niej
lepiej niż $\val{p15.partial.bound}$. To sprawdzenie arytmetyki, a nie idei,
ale właśnie ono łapie znak.

\begin{aibox}
To właśnie liczy framework dla każdego parametru modelu. Strata jest jedną
liczbą; każda waga jest jednym wejściem; a pochodna cząstkowa względem jednej
wagi jest odpowiedzią na \emph{jeśli poruszę tę jedną wagę i nic więcej, o ile
przesunie się strata}.

Jest ich tyle, ile wag, a ciekawym faktem o nich nie jest żadna pojedyncza
liczba, tylko to, że jeden przebieg wsteczny produkuje je wszystkie naraz
\dash{} co jest tematem Programu~\ref{prog:P16} i potrzebuje kształtu
obliczenia, a nie więcej rachunku różniczkowego.
\end{aibox}
\end{fr}

% ============================================================
\section{Zebranie ich w jeden obiekt}
\label{sec:P15-gradient}
\index{gradient}

\begin{fr}
Dwie pochodne cząstkowe w punkcie to dwie liczby. Zapisz je jako wektor, w
kolejności, w jakiej idą wejścia:
\[ \nabla f = \left(\frac{\partial f}{\partial x},\
                    \frac{\partial f}{\partial y}\right) \]
To jest \textbf{gradient}, a $\nabla$ czyta się \enquote{del} albo
\enquote{nabla}.

W punkcie z sekcji 1 wynosi on $(\val{p15.fx.at}, \val{p15.fy.at})$.
Program~\ref{prog:F09} dał wektorowi długość. Ile wynosi ta?
\dotline
\nextframe
\end{fr}

\begin{fr}
\ans{$\val{p15.grad.len}$}
Z $\sqrt{\val{p15.fx.at}^{2} + \val{p15.fy.at}^{2}}$.

Zatrzymaj fakt, że ma on długość oprócz kierunku, bo oba okażą się coś
znaczyć, a znaczenia są różne.

\begin{notationbox}
\textbf{Słowo \enquote{gradient} wykonuje w tej książce dwie prace i tutaj się
one spotykają.}

Program~\ref{prog:F06} i Program~\ref{prog:F11} używają go w zwykłym sensie:
\emph{współczynnik kierunkowy} prostej, jedna liczba, $m$ w $y = mx + c$. Tam
było to właściwe słowo i tamte programy nie są w błędzie.

Tutaj jest to \emph{wektor}, po jednej składowej na wejście. Oba są powiązane
\dash{} w jednym wymiarze gradient ma jedną składową i jest nią współczynnik
kierunkowy \dash{} ale nie są tym samym rodzajem obiektu, a zdanie w rodzaju
\enquote{gradient jest duży} znaczy co innego zależnie od tego, o który
chodzi. Gdy ten program mówi gradient, ma na myśli wektor.
\end{notationbox}
\end{fr}

\begin{fr}
Pytanie, które brzmi trywialnie i nie jest. Po co w ogóle zapisywać dwie
liczby jako wektor, zamiast trzymać je jako dwie liczby?
\dotline
\nextframe
\end{fr}

\begin{fr}
\begin{ansblock}
Bo wektor można pomnożyć skalarnie przez inny wektor, a dwóch luźnych liczb
nie można. Cała reszta tego programu wychodzi z tej jednej operacji.
\end{ansblock}

Program~\ref{prog:P05} poświęcił cały program temu, co kupuje iloczyn
skalarny, a odpowiedzią były długości, kąty i prostopadłość z jednej
definicji. Sekcja 3 wprowadza gradient do tej maszynerii i wydobywa z niej
najszybszy spadek, nie definiując niczego nowego.

\begin{note}
Przy więcej niż dwóch wejściach nie zmienia się nic poza długością listy.
Model z miliardem parametrów ma gradient o miliardzie składowych, nadal jest
to jeden wektor, a każde zdanie tego programu nadal jest o nim prawdziwe.

Książka zostaje w dwóch wymiarach, bo dwa to tyle, ile widać. Warto przy tym
pamiętać o ostrzeżeniu Programu~\ref{prog:P05}: dwuwymiarowe obrazy geometrii
wielowymiarowej mylą w konkretny, mierzalny sposób, a iloczyn skalarny jest
jedną z niewielu rzeczy, które przenoszą się bez zastrzeżeń.
\end{note}
\end{fr}

\begin{fr}
Jeszcze jedna własność, zanim zostanie użyty. Program~\ref{prog:F11}
powiedział, że pochodna stałej to zero. Ile wynosi gradient w punkcie, w
którym $f$ ma minimum?
\nextframe
\end{fr}

\begin{fr}
\begin{ansblock}
Wektor zerowy: każda pochodna cząstkowa jest tam zerem.
\end{ansblock}

Bo gdyby którakolwiek nie była, mógłbyś przesunąć się kawałek wzdłuż tego
wejścia i zejść w dół, więc nie byłbyś w minimum.

\begin{trapbox}
Odwrotne zdanie jest fałszywe i Program~\ref{prog:F11} powiedział to w jednym
wymiarze: zerowa pochodna to nie minimum. W więcej niż jednym jest gorzej, bo
istnieje kształt, który nie ma jednowymiarowego odpowiednika w ogóle \dash{}
punkt będący minimum wzdłuż jednego kierunku i maksimum wzdłuż innego.

Program~\ref{prog:P10} nazwał go i narysował: \emph{siodło}. Jego gradient
jest wektorem zerowym, więc cokolwiek zatrzymuje się, gdy gradient znika,
zatrzymuje się tam całkiem chętnie. Odróżnienie tych trzech wymaga drugich
pochodnych, co należy do Programu~\ref{prog:P17}.
\end{trapbox}
\end{fr}

% ============================================================
\section{Tempo zmiany wzdłuż dowolnego kierunku}
\label{sec:P15-directional}
\index{derivative!directional}

\begin{fr}
Dwie pochodne cząstkowe odpowiadają na dwa pytania: jak szybko $f$ się zmienia,
gdy idziesz wzdłuż $x$, i gdy idziesz wzdłuż $y$.

Żadne z nich nie jest pytaniem, które naprawdę masz. Chcesz poruszyć się w
jakimś wybranym przez siebie kierunku \dash{} powiedzmy na północny wschód albo
w którymkolwiek z nieskończenie wielu pozostałych \dash{} i wiedzieć, co robi
$f$.

Zapisz kierunek jako wektor jednostkowy $u$. Jak zdefiniowałbyś tempo zmiany
$f$ wzdłuż $u$, używając wyłącznie definicji pochodnej z
Programu~\ref{prog:F11}?
\dotline
\nextframe
\end{fr}

\begin{fr}
\begin{ansblock}
Dokładnie tak, jak zrobił to F11, tyle że wzdłuż tej prostej: zrób mały krok
$h$ w kierunku $u$, zobacz, o ile zmieniło się $f$, podziel przez $h$ i
zmniejszaj $h$.
\end{ansblock}

\[ D_u f = \lim_{h \to 0} \frac{f(\mathbf{p} + h u) - f(\mathbf{p})}{h} \]

To ta sama granica, na funkcji jednej zmiennej, którą dostajesz, idąc wzdłuż
prostej przez punkt. Nic nowego nie zostało zdefiniowane; wybrano kierunek i
zastosowano wzdłuż niego definicję F11.

\textbf{Jednostkowość ma znaczenie.} Gdyby $u$ było dwa razy dłuższe,
odpowiedź by się podwoiła i mówiłaby o długości twojej linijki, a nie o
stromości terenu.
\end{fr}

\begin{fr}
Teraz krok, dla którego ta sekcja istnieje. Ta granica jest kłopotliwa do
policzenia dla każdego kierunku, który mógłby cię obchodzić \dash{} i okazuje
się, że nigdy nie musisz.

Dwie pochodne cząstkowe to odpowiedzi dla $u = (1, 0)$ i $u = (0, 1)$. Ogólne
$u = (u_1, u_2)$ to $u_1$ pierwszego i $u_2$ drugiego. Ile wynosi więc
$D_u f$ w terminach gradientu?
\dotline
\nextframe
\end{fr}

\begin{fr}
\begin{ansblock}
\[ D_u f = \nabla f \cdot u \]
\end{ansblock}

Tempo wzdłuż dowolnego kierunku to gradient pomnożony skalarnie przez ten
kierunek, a dwie pochodne cząstkowe to przypadki szczególne $u = (1,0)$ i
$u = (0,1)$.

To cała ta sekcja i warto jasno powiedzieć, jakiego rodzaju jest to zdanie.
Nie jest definicją i nie jest konwencją: jest konsekwencją tego, że granica
jest liniowa względem kierunku, i można je sprawdzić. Skrypt bierze
$\val{p15.dir.steps}$ kierunków dookoła okręgu, liczy granicę numerycznie dla
każdego i porównuje z iloczynem skalarnym. Najgorsza niezgodność gdziekolwiek
jest poniżej $\val{p15.dir.bound}$.

\transcript{p15-directional}

\mermaidfig{p15-any-direction}{Jedna liczba na kierunek i skąd ona
pochodzi.}{dowolny kierunek}
\end{fr}

\begin{fr}
Przed konsekwencjami jeden rachunek. W tym samym punkcie, przy gradiencie
$(\val{p15.fx.at}, \val{p15.fy.at})$, ile wynosi tempo zmiany w kierunku
$u = (\num{0.6}, \num{0.8})$?

Sprawdź najpierw, czy $u$ jest wektorem jednostkowym, bo jeśli nie jest,
odpowiedź nic nie znaczy.
\dotline
\nextframe
\end{fr}

\begin{fr}
\begin{ansblock}
$\val{p15.fx.at} \times \num{0.6} + \val{p15.fy.at} \times \num{0.8} =
\num{17.6}$, a $u$ jest wektorem jednostkowym, bo
$\num{0.6}^{2} + \num{0.8}^{2} = 1$.
\end{ansblock}

Jest to liczba, którą drukuje listing powyżej, i jest mniejsza niż własna
długość gradientu, $\val{p15.grad.len}$. To nie przypadek i o tym jest
następna sekcja.

\begin{note}
Trójkąt $3$--$4$--$5$ jest powodem, dla którego $(\num{0.6}, \num{0.8})$ jest
wektorem jednostkowym, a Program~\ref{prog:P09} i Program~\ref{prog:P10}
budowały z tej samej trójki obroty bez zaokrągleń. To najtańszy dokładny
wektor jednostkowy, jaki istnieje, i dlatego wciąż się pojawia.
\end{note}
\end{fr}

% ============================================================
\section{Najbardziej stromo i w poprzek poziomicy}
\label{sec:P15-steepest}
\index{gradient!steepest direction}

\begin{fr}
Przeczytaj teraz $D_u f = \nabla f \cdot u$, mając przed sobą tożsamość z
Programu~\ref{prog:P05}:
\[ a \cdot b = \lVert a \rVert \, \lVert b \rVert \cos\theta \]
Skoro $u$ jest wektorem jednostkowym, $\lVert u \rVert = 1$, więc
\[ D_u f = \lVert \nabla f \rVert \cos\theta \]
gdzie $\theta$ to kąt między wybranym przez ciebie kierunkiem a gradientem.

Popatrz, co to zrobiło z pytaniem. Pierwszy czynnik jest ten sam, jakikolwiek
kierunek wybierzesz \dash{} jest własnością punktu, a nie twojego wyboru
\dash{} więc jedyne, czym twój wybór steruje, to $\cos\theta$. Całe
\enquote{który z nieskończenie wielu kierunków} zwinęło się w
\enquote{który cosinus}, a to jest pytanie o jednosłownej odpowiedzi.

Które $\theta$ czyni tempo największym i ile wynosi to największe tempo?
\dotline
\nextframe
\end{fr}

\begin{fr}
\begin{ansblock}
$\theta = 0$, a największe tempo to samo $\lVert \nabla f \rVert$.
\end{ansblock}

Bo $\cos\theta$ wynosi co najwyżej $1$ i równa się $1$ tylko wtedy, gdy oba
wskazują tak samo.

\textbf{Gradient jest więc kierunkiem najszybszego wzrostu, i jest to
wyprowadzone, a nie orzeczone.} Nie jest to definicja i nie jest to fakt do
przyjęcia na wiarę: wynika z tego, że tempo jest iloczynem skalarnym, i z
tego, że cosinus jest największy w zerze. To zdanie, na które zapracował cały
ten program.

A długość też teraz coś znaczy: $\lVert \nabla f \rVert$ mówi, \emph{jak
szybko} $f$ wspina się w tym najlepszym kierunku. Kierunek i długość, dwa
znaczenia, jeden wektor.
\end{fr}

\begin{fr}
Pomiar, bo \enquote{największe} to twierdzenie o każdym kierunku, a nie o
jednym. Wśród $\val{p15.dir.steps}$ kierunków największe znalezione tempo jest
tym pod $\val{p15.grad.deg}$ stopniami, czyli we własnym kierunku gradientu, a
jego wartość to własna długość gradientu.

Próbkowane maksimum ustępuje $\lVert \nabla f \rVert$ o mniej niż
$\val{p15.sweep.shortfall.bound}$, a skrypt orzeka, że ten niedobór nie
przekracza tego, co da się wyjaśnić rozdzielczością siatki \dash{} a nie, że
jest zerem, co byłoby fałszem dla każdego skończonego przeglądu.

\begin{note}
Pierwsza wersja tego sprawdzenia orzekała, że próbkowane maksimum
\emph{równa się} długości, i zawiodła. Zasłużyła na to: maksimum po trzech i
pół tysiąca kierunków nie może równać się maksimum po wszystkich, więc
tolerancja była liczbą dobraną tak, by sprawdzenie przeszło, a nie zdaniem o
matematyce.

Orzekana jest teraz para \dash{} że żaden kierunek nie bije gradientu, co jest
dokładne i jest nierównością z Programu~\ref{prog:P05}; oraz że próbkowany
najlepszy ustępuje o nie więcej niż rozdzielczość. Drugie zamienia niedobór z
zawstydzenia w wyjaśnienie.
\end{note}
\end{fr}

\begin{fr}
Teraz drugi koniec tego samego równania. Które $\theta$ czyni tempo zerem i co
to znaczy o tym, gdzie możesz iść?
\dotline
\nextframe
\end{fr}

\begin{fr}
\begin{ansblock}
$\theta = 90$ stopni. Idź pod kątem prostym do gradientu, a $f$ nie zmienia
się wcale, w pierwszym rzędzie.
\end{ansblock}

Droga, wzdłuż której $f$ się nie zmienia, to \textbf{poziomica} \dash{} zbiór
poziomicowy, linie na mapie, pierścienie wokół dna niecki. Zatem:

\textbf{gradient jest prostopadły do poziomicy przechodzącej przez punkt.}

To to samo równanie odczytane na drugim krańcu i warto zauważyć, że
\emph{najbardziej stromo pod górę} i \emph{prostopadle do poziomicy} nie są
dwoma faktami o gradiencie. Są jednym faktem, odczytanym przy
$\cos\theta = 1$ i przy $\cos\theta = 0$.
\end{fr}

\begin{fr}
Twierdzenie to można postawić dokładnie zamiast numerycznie i skrypt to robi.
Dla $f = ax^{2} + by^{2}$ gradient wynosi $(2ax, 2by)$, a kierunek wzdłuż
poziomicy to $(-by, ax)$. Pomnóż je skalarnie.
\nextframe
\end{fr}

\begin{fr}
\begin{ansblock}
$-2abxy + 2abxy = 0$. Nie w przybliżeniu zero \dash{} tożsamościowo zero, dla
każdego $a$, $b$, $x$ i $y$.
\end{ansblock}

Skrypt sprawdza to nad liczbami wymiernymi w $\val{p15.perp.trials}$
kombinacjach, bez żadnej tolerancji, bo twierdzenie zawierające słowo
\emph{prostopadły} jest twierdzeniem o dokładnym zerze, a porównanie
zmiennoprzecinkowe zamieniłoby je w twierdzenie o progu.

\begin{note}
To ta sama dyscyplina, której Program~\ref{prog:P04} użył dla rzędu, a
Program~\ref{prog:P13} dla rozkładu stacjonarnego: tam, gdzie matematyka jest
dokładna, licz dokładnie i orzekaj równość. Tolerancja jest dla pomiaru. To
nim nie jest.
\end{note}
\end{fr}

\begin{fr}
Oto dlaczego zaraz będzie to miało znaczenie, i warto odpowiedzieć, zanim
pójdziesz dalej.

Stoisz na powierzchni w kształcie niecki, gdzieś z boku. Gradient pod twoimi
stopami wskazuje pod kątem prostym do poziomicy, na której stoisz. Czy
kierunek prosto w dół tego zbocza wskazuje najniższy punkt niecki?
\dotline
\nextframe
\end{fr}

\begin{fr}
\begin{ansblock}
Tylko jeśli poziomice są okręgami. Ogólnie nie.
\end{ansblock}

Większość ludzi odpowiada tak i powód jest dobry: to prawda dla niecki, którą
każdy sobie wyobraża. Okrągła niecka ma poziomice będące okręgami, prostopadła
do okręgu wskazuje jego środek, a oba kierunki pokrywają się dokładnie.

\textbf{Ale prostopadłość do poziomicy jest zdaniem lokalnym, a to, gdzie jest
minimum, zdaniem globalnym}, i nic nie sprawia, że się zgadzają. Rozciągnij
nieckę wzdłuż jednej osi, a poziomice stają się elipsami; prostopadła do
elipsy nie wskazuje jej środka, chyba że akurat stoisz na jednej z jej osi.

Sekcja 5 mówi, co to robi z marszem.

\mermaidfig{p15-across-the-valley}{Gdzie wskazuje prostopadła, w okrągłej i w
długiej niecce.}{w poprzek doliny}
\end{fr}

% ============================================================
\section{Co to robi z marszem}
\label{sec:P15-zigzag}
\index{gradient!descent}

\begin{fr}
Program~\ref{prog:F11} zapisał już marsz, w jednym wymiarze:
\[ x \;\longleftarrow\; x - \eta\, f'(x) \]
i nazwał minus całym mechanizmem sterowania. W więcej niż jednym wymiarze
zmienia się tylko to, że pochodna jest wektorem:
\[ \mathbf{p} \;\longleftarrow\; \mathbf{p} - \eta\, \nabla f(\mathbf{p}) \]

Zauważ, że aktualizacja jest \emph{po współrzędnych}: każda współrzędna
przesuwa się o własną pochodną cząstkową i nie radzi się żadnej innej.
Dlatego framework może zastosować to do miliarda parametrów, nie musząc niczego
koordynować \dash{} notacja wektorowa jest wygodą przy zapisie, a arytmetyka
pod spodem to miliard niezależnych odejmowań.

Powiedz, dlaczego jest tam minus, w terminach wyniku tej sekcji, a nie wyniku
F11.
\dotline
\nextframe
\end{fr}

\begin{fr}
\begin{ansblock}
Bo gradient jest kierunkiem najszybszego \emph{wzrostu}, a ty chcesz iść w
dół.
\end{ansblock}

Jest to zdanie mocniejsze, niż F11 mógł postawić. F11 wiedział, że przeciwnie
do znaku pochodnej jest w dół; nie wiedział, że ujemny gradient jest
\emph{najlepszym} dostępnym kierunkiem w dół, bo to wymaga iloczynu skalarnego
i cosinusa. Teraz nie jest to reguła kciuka, tylko odpowiedź na pytanie,
\emph{który ze wszystkich kierunków schodzi najszybciej}.

\begin{trapbox}
\textbf{\enquote{Spadek gradientu porusza się w kierunku gradientu.}} Porusza
się w kierunku gradientu \emph{ujemnego}, a jest to najczęstszy błąd znaku w
ręcznie pisanej pętli treningowej.

Warto zobaczyć, ile to kosztuje, zamiast dać sobie to powiedzieć. W punkcie
$(\val{p15.zig.start.x}, \val{p15.zig.start.y})$ na niecce poniżej wartość
wynosi $\val{p15.sign.before}$. Jeden mały krok przeciwnie do gradientu
sprowadza ją do $\val{p15.sign.down}$. Jeden mały krok \emph{zgodnie} z
gradientem podnosi ją do $\val{p15.sign.up}$.

Nic nie zgłasza błędu. Pętla działa, liczby są skończone, a strata rośnie
\dash{} co czyta się jako zły krok uczenia albo zła inicjalizacja znacznie
częściej niż jako brakujący minus.
\end{trapbox}
\end{fr}

\begin{fr}
Teraz niecka. Program~\ref{prog:P10} zbudował jedną i celowo jej nie wydał:
funkcję kwadratową, której dwie wartości własne to $\val{p15.lam.hi}$ i
$\val{p15.lam.lo}$, więc jej poziomice są elipsami
$\val{p10.basin.axisratio}$ raza dłuższymi niż szerszymi. To dolina: stroma w
poprzek, płytka wzdłuż.

Ten program używa tej samej niecki, zamiast wymyślać kolejną, a skrypt czyta
zatwierdzone liczby P10 i przerywa budowanie, jeśli się przesunęły.

Zacznij w $(\val{p15.zig.start.x}, \val{p15.zig.start.y})$, czyli poza obiema
osiami. Zanim cokolwiek policzysz: czy ujemny gradient wskazuje tam minimum?
\nextframe
\end{fr}

\begin{fr}
\ans{Nie \dash{} jest $\val{p15.zig.angle}$ stopni od niego}
Minimum jest w początku układu, więc droga do domu prowadzi prosto z powrotem
wzdłuż prostej, na której stoisz. Ujemny gradient stoi za to prostopadle do
poziomicy, a na elipsie są to dwa różne kierunki.

\textbf{Czterdzieści dwa stopnie to nie błąd zaokrąglenia, to prawie połowa
kąta prostego.} Sterowanie nie jest lekko nietrafione; wskazuje zupełnie gdzie
indziej, i robi to z powodu, który jest teraz całkowicie zrozumiały, a nie
tajemniczy.

\begin{note}
Okrągła niecka jest próbą kontrolną i jest dokładna, a nie prawie dokładna.
Ustaw obie wartości własne równe, a gradient staje się wielokrotnością samego
położenia, więc ujemny gradient wskazuje minimum \emph{dokładnie}, a marsz
jest prostą linią. Skrypt orzeka to iloczynem wektorowym na floatach, który
jest dokładnie zerem.

Zapytanie arcus cosinusa o kąt tego nie mówi: raportuje około
$\val{p15.round.acos.bound}$ stopnia zamiast zera, bo cosinus wraca jako jeden
minus zaokrąglenie, a arcus cosinus liczby tak bliskiej $1$ to miejsce, gdzie
mieszka kasowanie z Programu~\ref{prog:P02}. Geometria jest dokładna; arcus
cosinus jest złym przyrządem do mierzenia zera.
\end{note}
\end{fr}

\begin{fr}
Zrób więc marsz. $\val{p15.zig.steps}$ kroków o rozmiarze
$\val{p15.zig.eta}$ z tamtego rogu, na tej niecce.

Dwie rzeczy do przewidzenia, zanim pójdziesz dalej. Ile razy marsz przechodzi
z jednej strony doliny na drugą? I jak dystans, który przebywa, ma się do
dystansu, który faktycznie pokonuje?
\dotline
\nextframe
\end{fr}

\begin{fr}
\begin{ansblock}
Przechodzi $\val{p15.zig.crossings}$ razy w $\val{p15.zig.steps}$ krokach
\dash{} co każdy krok \dash{} i przebywa $\val{p15.zig.ratio}$ raza dalej,
niż pokonuje.
\end{ansblock}

$\val{p15.zig.pathlen}$ drogi, żeby pokonać $\val{p15.zig.moved}$ terenu.
Prawie siedem ósmych marszu idzie w bok.

\begin{note}
Miarą jest dystans przebyty wobec dystansu \emph{pokonanego}, a nie wobec
dystansu do minimum. Wcześniejszy szkic użył drugiego i podał liczbę poniżej
jedynki dla marszu będącego prostą linią \dash{} bo po dwudziestu krokach nie
dotarł on do celu, więc przebył mniej niż całą drogę do domu. To mówi coś o
rozmiarze kroku, a nic o kierunku.

Droga przez przemieszczenie wynosi $1$ dla każdej prostej drogi, jak daleko by
nie zaszła, co jest dokładnie tą własnością, o której jest ta sekcja, i wynosi
$1$ dla okrągłej niecki co do ostatniego bitu.
\end{note}
\end{fr}

\begin{fr}
Mechanizm to jedna linia arytmetyki i jest to własna rekurencja
Programu~\ref{prog:F11}, po jednym kierunku naraz.

Na tej niecce oba kierunki nie oddziałują: każda współrzędna jest mnożona
przez $1 - \eta\lambda$ w każdym kroku, z własną $\lambda$. Przy
$\eta = \val{p15.zig.eta}$ policz ten czynnik dla kierunku stromego
($\lambda = \val{p15.lam.hi}$) i dla płytkiego
($\lambda = \val{p15.lam.lo}$).
\dotline
\nextframe
\end{fr}

\begin{fr}
\ans{$\val{p15.fac.hi}$ oraz $\val{p15.fac.lo}$}
I oto zygzak, w znakach. Stroma współrzędna mnożona jest przez liczbę
\emph{ujemną}, więc zmienia stronę w każdym kroku, kurcząc się o piątą część.
Płytka mnożona jest przez liczbę dodatnią tuż poniżej jedynki, więc pełznie ku
zeru, nigdy nie przechodząc na drugą stronę.

\textbf{Jeden kierunek przestrzeliwuje, a drugi się czołga, i odpowiada za oba
ten sam rozmiar kroku.} Zrób krok dość mały, by przestrzeliwanie ustało, a
płytki kierunek będzie szedł wiecznie; zrób go dość duży, by ruszyć płytki
kierunek, a stromy rozleci się.

\begin{rigourbox}
Ile faktycznie wynosi największe użyteczne $\eta$ i jak zależy od krzywizny,
należy do Programu~\ref{prog:P17}: potrzeba do tego porządnie drugiej
pochodnej, która jest macierzą i nie została zdefiniowana. Ten program może
powiedzieć, że czynnik staje się ujemny i że marsz przechodzi na drugą stronę;
nie może powiedzieć, gdzie leży granica.

A naprawa należy do Programu~\ref{prog:P20}. Moment pędu uśrednia ostatnie
kroki, więc składowe boczne \dash{} które zmieniają znak \dash{} kasują się
wzajemnie, podczas gdy te do przodu \dash{} które nie zmieniają \dash{}
sumują się. To zdanie staje się oczywiste dopiero po zobaczeniu tego obrazu i
dlatego ten program produkuje problem i zostawia go stojącym.
\end{rigourbox}
\end{fr}

\begin{fr}
Ostatnia rzecz, bo to powód, dla którego cokolwiek z tego się uogólnia.

Wszystko w tej sekcji zrobiono na dwóch wejściach i niecce, którą da się
narysować. Prawdziwa strata ma miliard wejść i nie da się jej narysować
wcale. Które z tych zdań to przeżywa, a które było o obrazku?
\dotline
\nextframe
\end{fr}

\begin{fr}
\begin{ansblock}
Wszystkie przeżywają. Ani jeden krok nie skorzystał z tego, że wejść były
dwa.
\end{ansblock}

Pochodna kierunkowa jest iloczynem skalarnym w dowolnej liczbie wymiarów;
cosinus jest największy w zerze w dowolnej; poziomica staje się powierzchnią, a
gradient nadal stoi do niej prostopadle; a dolina staje się kierunkiem, w
którym krzywizna jest mała w porównaniu z inną, a takich jest w bród, gdy do
wyboru jest ich miliard.

\textbf{Nie przeżywa obrazek, a nigdy nie był argumentem.} Rysunek był
sposobem zobaczenia iloczynu skalarnego. To iloczyn skalarny jest prawdą.

\begin{aibox}
Dlatego \enquote{krajobraz straty} jest zwrotem użytecznym i niebezpiecznym.
Sterowanie naprawdę jest tym, co opisał ten program, w dowolnie wielu
wymiarach. \emph{Kształt} nie: Program~\ref{prog:P05} zmierzył, jak źle
przenosi się dwuwymiarowa intuicja, a powierzchnia miliardowymiarowa ma
miejsce na zachowania nieposiadające dwuwymiarowego odpowiednika w ogóle.

Ufaj tu algebrze i nie ufaj obrazowi w głowie, w tej kolejności.
\end{aibox}
\end{fr}

% ============================================================
\begin{summarybox}
\sumitem{1--2}{\textbf{Pochodna cząstkowa trzyma każde inne wejście w miejscu}
  i jest wtedy zwykłą pochodną jednej zmiennej, więc każda reguła z
  Programu~\ref{prog:F12} stosuje się bez zmian.}
\sumitem{3--5}{Dla $f = x^{2}y + 3y$ pochodne cząstkowe to $2xy$ oraz
  $x^{2}+3$. \textbf{Trzymanie wejścia w miejscu nie usuwa go z odpowiedzi}
  \dash{} pochodna cząstkowa nadal jest funkcją wszystkiego.}
\sumitem{6--7}{W $(\val{p15.pt.x}, \val{p15.pt.y})$ wynoszą one
  $\val{p15.fx.at}$ i $\val{p15.fy.at}$, sprawdzone wobec różnicy centralnej
  na siatce lepiej niż $\val{p15.partial.bound}$.}
\sumitem{8--9}{\textbf{Gradient to pochodne cząstkowe zapisane jako wektor},
  $\nabla f$, mający długość oprócz kierunku \dash{} tutaj
  $\val{p15.grad.len}$.}
\sumitem{9}{\textbf{Słowo wykonuje dwie prace.} Programy F06 i F11 używają
  słowa \enquote{gradient} dla współczynnika kierunkowego prostej, jednej
  liczby; tutaj jest to wektor po jednej składowej na wejście.}
\sumitem{10--11}{Zapisuje się go jako wektor \textbf{po to, by dało się go
  pomnożyć skalarnie}, i stąd bierze się reszta programu.}
\sumitem{12--13}{W minimum gradient jest wektorem zerowym, a \textbf{odwrotne
  zdanie zawodzi} \dash{} siodło też ma gradient zerowy, a odróżnienie ich
  należy do Programu~\ref{prog:P17}.}
\sumitem{14--17}{\textbf{Tempo zmiany wzdłuż kierunku jednostkowego $u$ to
  $\nabla f \cdot u$}, czyli granica z Programu~\ref{prog:F11} wzięta wzdłuż
  prostej. Sprawdzone wobec tej granicy na $\val{p15.dir.steps}$ kierunkach
  lepiej niż $\val{p15.dir.bound}$.}
\sumitem{20--21}{Z cosinusem z Programu~\ref{prog:P05}
  $D_u f = \lVert\nabla f\rVert\cos\theta$, więc \textbf{kierunkiem
  najbardziej stromym jest własny kierunek gradientu, a najszybszym tempem
  jego długość}. Wyprowadzone, nie orzeczone.}
\sumitem{23--24}{Przy $\cos\theta = 0$ tempo jest zerem, więc \textbf{gradient
  jest prostopadły do poziomicy} \dash{} to samo równanie odczytane na drugim
  krańcu, dokładnie zero na $\val{p15.perp.trials}$ przypadkach wymiernych.}
\sumitem{25--26}{Prostopadłość jest sprawdzona \textbf{dokładnie nad liczbami
  wymiernymi}, bez żadnej tolerancji, bo twierdzenie zawierające słowo
  \emph{prostopadły} jest twierdzeniem o dokładnym zerze.}
\sumitem{27--28}{\textbf{Prostopadłość do poziomicy to nie to samo co
  wskazywanie minimum}, chyba że poziomice są okręgami. Jedno jest lokalne, a
  drugie globalne.}
\sumitem{29--30}{Marsz to marsz Programu~\ref{prog:F11} z pochodną
  wektorową, a minus ma teraz mocniejszy powód: gradient jest najszybszym
  \emph{wzrostem}.}
\sumitem{30}{\textbf{Znak jest najczęstszym błędem w ręcznie pisanej pętli.}
  Jeden krok we właściwą stronę sprowadza $\val{p15.sign.before}$ do
  $\val{p15.sign.down}$; jeden krok w złą podnosi to do
  $\val{p15.sign.up}$, bez żadnego zgłoszonego błędu.}
\sumitem{31--32}{Na niecce Programu~\ref{prog:P10}, o wartościach własnych
  $\val{p15.lam.hi}$ i $\val{p15.lam.lo}$, ujemny gradient wskazuje
  $\val{p15.zig.angle}$ stopni od minimum. W okrągłej niecce wskazuje je
  dokładnie.}
\sumitem{33--34}{Marsz przechodzi w poprzek doliny $\val{p15.zig.crossings}$
  razy w $\val{p15.zig.steps}$ krokach i \textbf{przebywa
  $\val{p15.zig.ratio}$ raza dalej, niż pokonuje}.}
\sumitem{35--36}{Każda współrzędna mnożona jest przez $1 - \eta\lambda$:
  czynnik wynosi $\val{p15.fac.hi}$ w kierunku stromym, który zmienia znak w
  każdym kroku, i $\val{p15.fac.lo}$ w płytkim, który pełznie. \textbf{Jeden
  rozmiar kroku musi obsłużyć oba.}}
\sumitem{37--38}{\textbf{Ani jeden krok argumentu nie skorzystał z tego, że
  wejść były dwa.} Obrazek jest ilustracją iloczynu skalarnego, a nigdy jego
  przesłanką.}
\end{summarybox}

\canyou

\begin{testexercises}
\item $g(x, y) = 3x^{2}y - y^{3}$. Zapisz obie pochodne cząstkowe i wylicz
  gradient w $(1, 2)$.
  \answerto{$\partial g/\partial x = 6xy$ oraz $\partial g/\partial y =
  3x^{2} - 3y^{2}$. W $(1, 2)$ daje to
  \[ \nabla g = (12,\ -9), \qquad \lVert \nabla g \rVert = 15. \]
  Znowu trójkąt $3$--$4$--$5$, przeskalowany przez trzy.}
\item W punkcie, w którym $\nabla f = (3, 4)$, ile wynosi tempo zmiany $f$ w
  kierunku $(\num{0.8}, \num{0.6})$, a ile w kierunku
  $(-\num{0.8}, -\num{0.6})$?
  \answerto{$\num{4.8}$ oraz $-\num{4.8}$. Odwrócenie kierunku odwraca znak i
  nic więcej, bo iloczyn skalarny jest liniowy względem $u$ \dash{} i dlatego
  też kierunek najszybszego \emph{spadku} jest dokładnie przeciwny do
  kierunku najszybszego wzrostu, a nie tylko blisko niego.}
\item Kolega mówi, że jego strata spada, ale bardzo wolno, i że długość
  gradientu jest duża. Co możesz mu powiedzieć na podstawie samego tego
  programu?
  \answerto{Że duży gradient i wolny postęp nie są sprzecznością, a prawdopodobnym
  powodem jest to, że większość gradientu wskazuje w poprzek doliny, a nie
  wzdłuż niej. Długość mówi, jak stromy jest teren w kierunku najbardziej
  stromym; nie mówi nic o tym, czy ten kierunek prowadzi dokądkolwiek. Co z
  tym zrobić, należy do Programu~\ref{prog:P20}.}
\item Dlaczego porównywanie $D_u f$ dla dwóch kierunków jest błędne, jeśli
  jeden z nich nie jest wektorem jednostkowym?
  \answerto{Bo $D_u f$ skaluje się z długością $u$, więc dłuższe $u$ podaje
  większą liczbę bez tego, żeby teren był choć trochę bardziej stromy.
  Porównywałbyś linijki, a nie zbocza. Wymóg jednostkowości jest tym, co czyni
  tę liczbę własnością powierzchni.}
\item W niecce o poziomicach będących okręgami ile razy spadek gradientu
  przechodzi z jednej strony na drugą, zanim dotrze?
  \answerto{Ani razu, dla dowolnego zbieżnego rozmiaru kroku. Ujemny gradient
  wskazuje środek dokładnie, więc każdy krok idzie wzdłuż tej samej prostej, a
  droga wynosi $1$ raza własne przemieszczenie \dash{} co skrypt orzeka
  dokładnie. Przechodzenie wymaga poziomic nieokrągłych i to cała treść
  sekcji 5.}
\end{testexercises}

\begin{furtherproblems}
\item Pokaż, że jeśli $u$ jest wektorem jednostkowym, to
  $D_{-u} f = -D_u f$, i powiedz, co to znaczy o kierunku najszybszego
  spadku.
  \answerto{$\nabla f \cdot (-u) = -(\nabla f \cdot u)$, bo iloczyn skalarny
  jest liniowy względem drugiego argumentu. Kierunkiem najszybszego
  \emph{spadku} jest więc dokładnie $-\nabla f$: ten sam argument, który
  uczynił $\cos\theta = 1$ przypadkiem najlepszym, czyni $\cos\theta = -1$
  najgorszym, a najgorsze tempo wzrostu jest najlepszym tempem spadku.}
\item Gradient $f(x,y) = ax^{2} + by^{2}$ w $(x, y)$ wynosi $(2ax, 2by)$. Dla
  których punktów jest on równoległy do kierunku z powrotem do początku
  układu?
  \answerto{Tam, gdzie $2ax \cdot y = 2by \cdot x$, czyli $xy(a - b) = 0$.
  Albo więc $a = b$ \dash{} okrągła niecka, gdzie jest równoległy wszędzie
  \dash{} albo punkt leży na jednej z dwóch osi. Wszędzie indziej na niecce
  wydłużonej oba kierunki się różnią, i dlatego marsz zaczynający poza osiami
  zygzakuje, a zaczynający na osi nie.}
\item Funkcja trzech wejść ma gradient o trzech składowych. Ile kierunków
  stoi do niego pod kątem prostym i co to znaczy o poziomicy?
  \answerto{Cała płaszczyzna ich, a nie dwa odosobnione kierunki: w trzech
  wymiarach wektory prostopadłe do danego tworzą przestrzeń dwuwymiarową, co
  jest liczeniem wymiaru z Programu~\ref{prog:P04}. Zbiór poziomicowy przez
  punkt jest więc \emph{powierzchnią}, a nie krzywą, a gradient stoi
  prostopadle do całej naraz.}
\item Dlaczego argument z sekcji 4 nigdy nie wspomina, ile wejść ma $f$?
  \answerto{Bo używa tylko tego, że $D_u f$ jest iloczynem skalarnym i że
  iloczyn skalarny to $\lVert a\rVert\lVert b\rVert\cos\theta$, a żadne z tych
  zdań nie ma w sobie wymiaru. To właśnie czyni wynik użytecznym dla modelu o
  miliardzie parametrów, gdzie nic nie da się narysować. Obrazek w sekcji 5
  jest ilustracją algebry i nigdy jej przesłanką.}
\item Masz nieckę o wartościach własnych $\lambda$ i $1$ i podwajasz
  $\lambda$. Co dzieje się z kątem między ujemnym gradientem a drogą do domu i
  jak byś to sprawdził?
  \answerto{Rośnie: im bardziej wydłużona niecka, tym dalej od środka wskazuje
  prostopadła. Skrypt sprawdza dokładnie to, przeglądając stosunek i orzekając,
  że kąt rośnie na każdym kroku \dash{} co jest niezmiennikiem, gdzie
  pojedynczy kąt byłby faktem o jednej niecce. Nie może urosnąć powyżej kąta
  prostego, bo krok szedłby wtedy pod górę.}
\end{furtherproblems}
